% From Message Authorization to Effect Binding — arXiv preprint
% Build: tectonic main.tex  (XeTeX engine; fetches packages automatically)
% STATUS: working draft circulated for co-author review and consent prior to submission.
%
% NOTE ON THE PREAMBLE. This mirrors ../../poc-research-roadmap/ov-poc-standard/paper/main.tex
% with one deliberate omission: `minted` is not loaded. It shells out to
% pygments, which would make the CI paper build depend on a Python package
% nobody installs there. `listings` needs no external tool and this paper has
% no syntax-highlighted listings to lose.
\documentclass[11pt]{article}

\usepackage[margin=1in]{geometry}
\usepackage{tikz}
\usepackage{graphicx}
\usepackage{booktabs}
\usepackage{array}
\usepackage{amsmath, amssymb}
\usepackage{listings}
\usepackage{xcolor}
\usepackage{colortbl}
\usepackage[hidelinks]{hyperref}
\usepackage{enumitem}
\usepackage{microtype}
\usepackage{caption}
\usepackage{siunitx}
\usepackage{amsthm}
\newtheorem{definition}{Definition}
\newtheorem{theorem}{Theorem}
\newtheorem{proposition}{Proposition}
\usepackage[section]{placeins}

% Advanced AI Society brand palette
\definecolor{aailime}{HTML}{CFFF04}
\definecolor{aailimetint}{HTML}{F2FFB8}
\definecolor{aaiink}{HTML}{0A0A0A}
\definecolor{aaired}{HTML}{FFE3E4}
\definecolor{aaineutral}{HTML}{E7E3DD}
\definecolor{aaipaper}{HTML}{F0EDEA}
\definecolor{aaimuted}{HTML}{8F8A82}
\definecolor{aaidarklime}{HTML}{7A9900}
\definecolor{codepanel}{HTML}{F7F5F1}

\lstset{basicstyle=\ttfamily\footnotesize, breaklines=true, frame=single,
        framerule=0.3pt, xleftmargin=1em, columns=fullflexible}
\captionsetup{font=small, labelfont=bf}

\newcommand{\transit}{\texttt{transit}}

% The Society's bracket mark, drawn natively (lime on the dark band)
\newcommand{\aaimark}[1]{%
\begin{tikzpicture}[x=#1, y=#1, line cap=butt]
  \fill[aailime] (0,0) rectangle (0.16,1);
  \fill[aailime] (0.16,0) rectangle (0.42,0.24);
  \fill[aailime] (0.16,0.76) rectangle (0.42,1);
  \fill[aailime] (1.04,0) rectangle (1.20,1);
  \fill[aailime] (0.78,0) rectangle (1.04,0.24);
  \fill[aailime] (0.78,0.76) rectangle (1.04,1);
\end{tikzpicture}}

\begin{document}

% ---------------------------------------------------------------- cover page
\thispagestyle{empty}

\begin{tikzpicture}[remember picture, overlay]
  \fill[aaiink] (current page.north west) rectangle
    ([yshift=-6.35cm]current page.north east);
  \fill[aailime] ([yshift=-6.35cm]current page.north west) rectangle
    ([yshift=-6.53cm]current page.north east);
  \node[anchor=north west, inner sep=0pt]
    at ([xshift=1in, yshift=-0.95cm]current page.north west) {\aaimark{0.62cm}};
  \node[anchor=north west, inner sep=0pt, text width=14.2cm]
    at ([xshift=1in+1.15cm, yshift=-0.98cm]current page.north west)
    {\sffamily\bfseries\footnotesize\color{aailime}%
     ADVANCED AI SOCIETY\; \textcolor{aaimuted}{$\cdot$}\; PROOF-OF-CONTROL
     INITIATIVE\\[3pt]
     {\normalfont\sffamily\footnotesize\color{aaimuted}
      Working draft v0.1 $\cdot$ August 2026 $\cdot$ for co-author review
      prior to submission}};
  \node[anchor=north west, inner sep=0pt, text width=\textwidth]
    at ([xshift=1in, yshift=-2.55cm]current page.north west)
    {\sffamily\bfseries\fontsize{21.5}{26}\selectfont\color{aaipaper}%
     From Message Authorization to Effect Binding: Conditions,
     Counterexamples, and a Conformance Harness};
  \node[anchor=north west, inner sep=0pt, text width=\textwidth]
    at ([xshift=1in, yshift=-5.55cm]current page.north west)
    {\sffamily\small\color{aaipaper}%
     \textcolor{aailime}{\rule[0.5ex]{1.6em}{2pt}}\hspace{0.6em}%
     Signing the bytes is not signing the deed.};
\end{tikzpicture}

\vspace*{5.9cm}

% Compact author block
{\small\noindent
\textbf{Jim Schwoebel}\textsuperscript{1,2}

\smallskip\noindent
{\footnotesize
\textsuperscript{1}\,Quome \quad
\textsuperscript{2}\,Advanced AI Society, \url{https://advancedaisociety.org}\footnote{Prepared
within the Proof-of-Control Initiative of the Advanced AI Society. Authorship is
limited to contributors to this manuscript. Co-authors will be added only upon
explicit consent and substantive contribution; the author list is deliberately
short at this stage rather than inherited from a sibling paper.}}}

\vspace{0.9em}
\noindent\textcolor{aailime}{\rule{\textwidth}{1.6pt}}
\vspace{0.2em}

\begin{center}
{\sffamily\bfseries\small ABSTRACT}
\end{center}
\vspace{-0.5em}
\noindent{\small
Here is something that sounds like it must be true, and isn't. You give a
program permission to do a thing. It signs a message saying what it wants. The
message arrives with the signature intact, nobody tampered with anything. So
the thing that happens is the thing you approved---right?

No. And it fails even when nobody is attacking you. Take a network retry over
an operation that isn't idempotent. Retry with the same one-time nonce and the
endpoint tells your client the operation failed---a lie, because it already
committed. Retry with a fresh nonce and it does the whole thing twice. Both
branches are wrong and there is no third branch. Every component followed its
manual. You cannot fix it by improving the message, because on the wire an
attacker replaying you and a load balancer retrying you are the same event;
the difference is somebody's intent, and intent is not a field in a JSON
object.

So we ask what would have to be true for the approved message to determine the
effect, and give four conditions: put every input in the message, give one
meaning one spelling, make redelivery harmless, and name the state you are
assuming. Together they are enough. We are loud about not having shown them
necessary---the document we inherited claims both directions and established
one.

Then we built it and tried to break it. There is a JSON decoder you can dial
to behave like Go, Node, Python, Ruby or Elixir, checked against the real
runtimes rather than asserted. There is an endpoint built to be bad in five
separate ways you can switch on one at a time. And there is a proxy for real
APIs, which closes one condition outright, closes another only if you
configure it, and cannot close the third at all---not ``doesn't yet,'' cannot.
We wrote a test asserting that third attack \emph{succeeds}, so the claim
cannot quietly grow later.

Finally we recounted a 24-endpoint survey we had inherited and the table does
not say what the prose around it says: 4~/~7~/~13, not 4~/~6~/~14, one row
dropped from one bucket and absorbed into another. We report that rather than
quietly fixing it. We also say what the corrected number still does not
license: two dozen endpoints picked for industry coverage is not a random
sample, so 54\% is no more a fact about the world's APIs than 58\% was.}

\clearpage

\section{The Problem}
\label{sec:problem}

I want to start with a case where nobody does anything wrong. That is
important. It is easy to look at a security paper, find the villain, and
conclude that the lesson is ``don't be that guy.'' Here there is no villain.
Everybody follows the manual and the result is still broken, and that tells
you the trouble is in the arrangement rather than in anyone's carelessness.

So: a program holds a token that lets it place an order. The token is narrow,
it expires quickly, and it carries a number that may be used exactly
once---a nonce---so that nobody can record the request and send it again
later. The program signs the bytes of its request. The endpoint checks the
signature: good. It checks the nonce: never seen before, good. It writes the
order into the database. It starts composing a reply.

The connection drops.

Now something between the two machines---a load balancer, a service mesh, a
client library with a perfectly sensible default---decides to try again. And
here is the whole paper in one paragraph: it has two ways to try again, and
both of them are wrong.

Suppose it sends the identical request, same nonce and all. The endpoint looks
at the nonce, sees that it has already been used, and refuses. Now, that is
the correct behaviour! That is precisely what you built the nonce for. If
somebody had recorded your request and was replaying it, you would want
exactly this. But nobody recorded anything. The order is sitting in the
database. Your program has now been told, in a well-formed and correctly
signed message, that the operation failed. It didn't fail. It happened.

All right, suppose instead the program gets a fresh nonce and tries again. Now
the endpoint sees a request that is legitimate, correctly signed, and has
never been seen before---because it hasn't---so it does what you told it to
do and places the order. Two orders. One customer. One intention.

Look at what we have. The signature check: right. The nonce check: right, and
doing its job well. The retry: right, that's what reliable transports are
\emph{for}. Every piece is correct and the assembly is wrong. And you cannot
repair it by making the message fancier, because the message layer is blind
to the distinction that matters. To an attacker replaying you and to a load
balancer retrying you, the bytes on the wire are identical. They are the same
event. The difference lives in somebody's intention, and intention is not a
field you can add to a JSON object.

That is the whole shape of the thing. Let us give the property we want a name:
\emph{effect binding}, meaning that the message you approved determines the
effect that happens. What the retry case shows is that message
authorization---which is what capability protocols give you, and they give it
to you well---is a different property, and no amount of strengthening the
message will turn one into the other.

There is a repair, and it is worth showing now because every repair in this
paper has the same shape. The mistake is that we used one number for two jobs.
The nonce is supposed to answer ``is this request fresh?'' and we quietly
asked it to also answer ``is this the same order?'' Those are different
questions. So use two numbers. Let the nonce change on every retry, as it
should. Let the client also carry a key that names \emph{the order}, and hold
that key constant across retries. Then the fresh-nonce retry shows up carrying
an order key the endpoint has seen before, and the endpoint hands back the
result it computed the first time instead of doing the work again. ISO~20022
has had this for years---they call it the end-to-end
identifier~\cite{iso20022}---and HTTP has it as the
\texttt{Idempotency-Key} header.

That closes one of the two branches. It does not close the other, and I want
to be honest about this rather than let you discover it later. If the first
reply never arrives, your program still does not know whether the effect
happened. It can ask, of course. But asking is another message, and that
message can be lost too, and so can the answer. You can push the uncertainty
around but you cannot make it go away. What is left over, after every fix an
endpoint can make, is a committed operation that the client believes failed.

We report that leftover separately from everything else the tool finds. A
well-behaved endpoint shows no findings, and if we printed ``no findings'' and
stopped, we would be telling you the problem was gone. It isn't gone. So the
tool prints the leftover under its own heading, every time, even on a clean
run.

\subsection{Now suppose somebody is attacking you}

That was what goes wrong with nobody attacking. Here is what goes wrong when
somebody is, and it is more embarrassing.

Two programs both read JSON. One decides whether the request is allowed; the
other carries it out. In a real deployment they are usually written in
different languages, because different teams built them. Look at this:

\begin{lstlisting}
{"to":"acct-1","amount":100,"to":"acct-9"}
\end{lstlisting}

The key \texttt{to} is in there twice. Is that legal? RFC~8259 says yes and
declines to say which one wins. RFC~7493 says no, don't do
that~\cite{rfc8259,rfc7493}. So the standards disagree, which means the
implementations disagree, and---this is the bad part---they disagree without
complaining. Go's \texttt{encoding/json} keeps the last one it sees.
Erlang/OTP's \texttt{json} module, which is what Elixir uses, keeps the first
one. Neither is a bug. Both are reasonable ways to read a document whose own
specification shrugged.

So now: the program that decides policy reads ``pay \texttt{acct-1}'', thinks
that looks fine, and signs. The program that moves the money reads ``pay
\texttt{acct-9}'' and pays \texttt{acct-9}. The signature verifies perfectly.
Not one byte was altered anywhere. Both programs did exactly what their
documentation promised. And your money went to the wrong place.

People call this a parser differential~\cite{kaminsky2010}. The name doesn't
matter much. What matters is that it is an old family---HTTP request smuggling
is the same trick played on a front-end and a back-end that disagree about
where one request stops and the next starts~\cite{linhart2005,kettle2019}. What
is new is the arrangement. When a person clicks a button, the thing that shows
the button and the thing that handles the click are usually the same program,
so there is nothing to disagree with. When a program signs a request on your
behalf, the signer and the doer are almost never the same program. We have
taken a fifty-year-old class of bug and built an industry on top of the
conditions that produce it.

\section{What Others Have Built}
\label{sec:related}

None of this is new territory. Smart people have been working on the
surrounding problems for fifty years, and I want to go through what they built
and be careful about what each thing does and does not settle. Being careless
here is how you end up thinking a problem is solved because something nearby
is solved.

\paragraph{Complete mediation, and the assumption hiding inside it.}
Saltzer and Schroeder put complete mediation near the top of their list of
design principles: check every access to every object, every
time~\cite{saltzer1975}. The reference monitor that formalizes
this~\cite{anderson1972} carries an assumption so obvious in 1972 that nobody
wrote it down, which is that the monitor and the thing being protected are the
same machine. The monitor looks at the request; the machine performs it. There
is no gap between deciding and doing because there is no distance between them.

Now take that monitor and split it across a network. The decision happens in
one process, on one host, in one programming language. The effect happens
somewhere else, in a different language, written by different people. The gap
that the 1972 model defined out of existence is precisely the gap I have been
describing. Complete mediation isn't wrong. It simply has nothing to say about
the case where the decider and the doer are two programs that have to agree on
what a request means---and, as we have just seen, sometimes don't.

\paragraph{Capabilities.}
This runs from Dennis and Van Horn~\cite{dennis1966} through Levy's
survey~\cite{levy1984} to macaroons~\cite{birgisson2014}, which let a holder
narrow a token but never widen it. That is a genuinely good answer to the
question ``what is this bearer allowed to ask for?'' It is not an answer to
``what will the endpoint do with what it asked?'' The caveat language has no
words for effects. It could not have; that is not what it is for.

\paragraph{Proofs about messages, which are proofs about messages.}
The Dolev--Yao model~\cite{dolevyao1983} and the tools built on
it---ProVerif~\cite{blanchet2001}, Tamarin~\cite{meier2013}---are excellent.
I mean that. If you want to know what an attacker can do to your protocol
messages, these are the right instruments and they are far better than
intuition.

But I want to be careful, because a green checkmark from a good tool is easy
to overread. When Tamarin tells you a capability protocol is secure, it is
telling you something about messages: which ones an attacker can manufacture,
which secrets stay secret, who agrees with whom about what. The endpoint's
parser is not in the model. Its defaulting behaviour is not in the model. Its
database is not in the model. And the retry case from
Section~\ref{sec:problem} is completely invisible to such a proof, because
both branches are protocol-legal. Nothing was forged. The proof is correct and
the system is broken, and those two facts sit together without contradiction
because they are about different things.

\paragraph{The confused deputy.}
Hardy's compiler~\cite{hardy1988} was talked into using its own authority for
a caller who didn't have any. Same skeleton as ours: a component with power
acts on a request whose meaning it got wrong. The difference is which part it
got wrong. Hardy's deputy was confused about \emph{whose} authority was in
play. Ours is confused about \emph{what was asked}.

\paragraph{Canonical form.}
RFC~8785 defines a canonical way to write JSON down~\cite{rfc8785}, and
RFC~7493 defines the subset of JSON on which that is even
meaningful~\cite{rfc7493}. Together they are the best available answer to the
duplicate-key problem, and one of our practical conclusions is dull enough to
state in a sentence: decode strictly and canonicalize \emph{before} you check
the signature, not after. It is the cheapest fix in this paper and probably
the one with the largest effect.

\paragraph{RFC 9110.}
HTTP already says which of its methods are idempotent and which are
not~\cite{rfc9110}. \texttt{PUT} is; \texttt{POST} isn't. This is normative
text, not folklore, and we lean on it directly in condition C3.

\paragraph{ISO 20022, which mostly solved this in the 2000s.}
Of everything in this section, this~\cite{iso20022} is the strongest existing
answer, and it deserves reading rather than citing. Look at what the
\texttt{pain.001} credit transfer actually does. The schema is mandatory and
closed, so there are no optional fields for a server to helpfully fill in on
your behalf. It carries an end-to-end identifier that is separate from
anything the transport does about retries. And that identifier is required,
not optional, so a client cannot forget it.

That is three of our four conditions, met by a message format designed by
people who could not afford to get this wrong, roughly two decades before
anybody was worried about software agents. So the conditions are achievable.
That is worth knowing, and it is the main reason I do not think this paper is
describing an impossibility.

It is also worth noticing what ISO~20022 is not. It is not a general-purpose
web API. Nobody designs a REST endpoint that way on a Tuesday afternoon, and
the interesting question is how much of that discipline an ordinary API can
adopt without becoming a payments network.

\section{Writing It Down Precisely}
\label{sec:conditions}

Now I have to write down some symbols, and I want to say why before I do it.
The symbols are not going to make anything truer than it was in
Section~\ref{sec:problem}. Everything important has already been said in
English. What the notation buys us is that it makes it much harder to be
vague---in particular, it forces us to write down exactly which things the
effect is allowed to depend on, and that turns out to be the entire question.
If you find the algebra unhelpful, skip to the four conditions; you will lose
nothing but precision.

Think of an endpoint as a machine that takes a state and some bytes and gives
you back a new state and something to look at. Let $\Sigma$ be the space of
states the endpoint can be in, $W$ the byte strings that can arrive on the
wire, and $O$ the things you can observe coming back. Parsing turns bytes into
a document, and it is a \emph{partial} function $P : W \rightharpoonup D$
because plenty of byte strings are not documents at all. Then the endpoint is

\[
E \;:\; \Sigma \times W \;\to\; \Sigma \times O .
\]

Split that into two steps, $E(\sigma, w) = X(\sigma, P(w))$: first read the
bytes, then act on what you read. And here is the part that matters. The
program that decides policy has \emph{its own} parser, which I will call $P'$,
and a policy $\Pi$ that looks at documents and approves or refuses. It signs
$w$ when $\Pi(P'(w))$ says yes.

Two parsers. $P$ and $P'$. Everything in this paper is downstream of the fact
that we wrote two of them and then assumed they were one.

\begin{definition}[Message authorization]
A deployment provides message authorization if no party can cause the endpoint
to accept a wire representation $w$ that was not signed by an authorized
holder.
\end{definition}

\begin{definition}[Effect binding]
A deployment provides effect binding if, for every authorized $w$, the
resulting state transition and output $E(\sigma, w)$ are determined by
$P'(w)$---the document the policy evaluated---and by nothing else.
\end{definition}

These two are independent, and that independence is the entire paper. The
retry case has message authorization. Nobody forged anything; every signature
checked out. It does not have effect binding. If you only ever ask for the
first property you will get it, and you will believe you got the second one,
and nothing will tell you otherwise until the money is in the wrong account.

So what would you have to demand to get the second one? Four things. Each is
pinned to a function that checks it in \texttt{src/classify.rs}, so that if
you disagree with a condition you can go read the code that enforces it rather
than argue with a paragraph.

\begin{description}[leftmargin=1.6em, style=nextline]
\item[C1 --- Put every input in the message.]
Everything that changes what happens has to be in the document that got
signed. Written out: $X(\sigma, d)$ must not read anything out of $\sigma$
that the client could have said and didn't. The way this fails is so ordinary
that it feels like a feature. You leave out \texttt{fee\_tier}; the server
looks up your account preferences and fills in ``priority.'' The policy
approved an ordinary transfer. The endpoint performed an expedited one. Nobody
lied to anybody. Checked by \texttt{judge\_c1}.

\item[C2 --- One meaning, one spelling.]
$P$ and $P'$ have to agree, and each has to be injective on the meanings it
accepts. That is two separate demands. Two different byte strings must not
mean the same thing, and---this is the duplicate-key case---one byte string
must not mean two different things to two different readers. The first half is
violated by a resource you can reach as both \texttt{/v1/users/me} and
\texttt{/v1/users/usr\_99812}; those are one thing wearing two names, and a
signature over one of them says nothing about the other. Checked by
\texttt{judge\_c2}.

\item[C3 --- Make redelivery harmless, with two numbers not one.]
Deliver the same authorized document ten times, get one effect. And the
identifier that decides ``this is the same document'' must be a different
thing from the nonce that decides ``this is fresh.'' That second clause is the
whole repair from Section~\ref{sec:problem}, and leaving it out is what
produces an endpoint that is idempotent in the brochure and not in practice.
Checked by \texttt{judge\_c3}, which grants this for free to the methods
RFC~9110 already defines as idempotent, because there the specification has
done the work.

\item[C4 --- Say what state you are assuming.]
If the effect depends on something that can change between the moment you sign
and the moment the endpoint acts, then the document has to name the state it
is counting on, and the endpoint has to refuse when that has moved. In HTTP
you spell this \texttt{If-Match}. The point is not the header. The point is
that ``transfer \$100 if the balance is still what I saw'' is a different
instruction from ``transfer \$100,'' and only one of them can be honoured
safely by a machine acting minutes later. Checked by \texttt{judge\_c4}.
\end{description}

\section{Sufficient, and Not Known To Be Necessary}
\label{sec:sufficiency}

Four conditions is a satisfying number. It has the shape of a finished result,
and there is a real temptation, having got there, to say ``an endpoint is
effect-binding \emph{if and only if} C1 through C4 hold'' and go to lunch.

I want to resist that, because we have not earned the second half.

What we have is: if all four hold, you get effect binding. That is the
direction with the engineering value in it, and the argument is short. If
every input that matters is in the document (C1), and the document has exactly
one spelling and exactly one meaning (C2), and delivering it repeatedly cannot
multiply the effect (C3), and the state it assumes is named and checked (C4),
then what the endpoint does is a function of the document the policy looked
at. That is Definition~2. That is the whole proof, and its shortness is a
point in its favour rather than against it.

What we do not have is the converse: that an endpoint failing one of these
conditions must therefore fail to be effect-binding. And it is not merely
unproven, it is doubtful, and I can give you two counterexamples off the top
of my head.

Put the monitor on the same machine as the resource. Now you are back in 1972,
the reference-monitor assumption holds again, and there is no gap for any of
C1--C4 to close. The endpoint is effect-binding and satisfies none of our
message-level conditions, because it doesn't need to.

Or take an endpoint whose effect is a total function of its input over a fixed
state space---no stored state that can move. It satisfies Definition~2 without
pinning anything, because there is nothing to pin. C4 is vacuous there.

Neither of those is exotic. Someone will build both.

The document this work grew out of states the result as an ``if and only if''
while establishing one direction, and I want to name that plainly rather than
inherit it quietly. It is the kind of thing a good referee finds in the first
ten minutes, and it should not require a referee.

Does the distinction matter, or is this pedantry? It matters, and specifically
it matters for a standard rather than for a proof. Write C1--C4 into a
standard as \emph{necessary} conditions and you have built a conformance
regime that refuses to certify systems which are, in fact, doing the right
thing by other means. You would be telling an engineer that their co-located
monitor is insufficient. It is sufficient. Your standard is wrong.

So: sufficient set, conformance target, and the converse marked as open. If
somebody proves necessity later, good---and if somebody produces a third
counterexample instead, that is good too, and rather more likely.

\section{How We Checked}
\label{sec:method}

Conditions on paper are cheap. We wanted to find out whether they survive
contact with real software, so we built a harness called \transit{}. It has
three parts, and how they relate to each other \emph{is} the method---so I
want to go through not just what each one does but why it is built that way
and what would have gone wrong otherwise.

\paragraph{One parser with knobs on it.}
We could have written the differential test by requiring five language
runtimes to be installed. Then the test would pass or fail depending on what
happens to be on your laptop, which is a poor property for evidence.

Instead there is one JSON parser with the disagreements turned into settings:
what to do with a duplicate key, what to do with a lone surrogate, how big a
number may get, whether to normalize Unicode. Dial it one way and it behaves
like Go. Dial it another and it behaves like Erlang.

Now, that is a claim, not a fact. When I say ``this configuration behaves like
Go,'' I am asserting something about Go that I could be wrong about, and if I
am wrong the whole differential is theatre. So we try to prove ourselves
wrong. Where a runtime \emph{is} installed, we shell out to it, push the same
corpus through, and demand that the imitation and the real thing agree. Where
the runtime is absent the model still runs, but every row of output carries the
word \texttt{Unverified}, in every format, forever, until somebody actually
checks it.

Here is a case in point, and I would rather tell you than have you find it.
The famous example in this area is Go's last-key-wins against Elixir's
first-key-wins. There is no Elixir on our development machine. So the
acceptance test does \emph{not} assert that pair. It asserts a disagreement
between two decoders that really are present and really were written
independently---a strict I-JSON decoder and \texttt{serde\_json}---and it
exercises first-key-wins through the Elixir model, which prints
\texttt{Unverified} wherever it appears. We have since measured OTP~27
directly on a different machine, and it does keep the first key. That is a
fact about that machine. It does not promote the model on this one, and the
tool does not pretend otherwise.

This is fussier than it needs to be if you trust us. The point is that you
shouldn't have to.

\paragraph{An endpoint built to be bad.}
The thing under test is a function: state and request go in, new state and
response come out. No sockets in the interesting path, so the tests cannot
flake, and the same function is what gets served over HTTP when you want to
poke it by hand.

It has five separate ways of being wrong---sloppy decoding, filling in
defaults from your profile, treating two URLs as one resource, forgetting
idempotency, ignoring preconditions---and each one switches independently. Turn
them all off and you get what we call the \emph{boundable profile}.

That profile is the control, and controls are the part people skip. Suppose I
write a check for the defaulting bypass and it reports a finding. Have I
learned something about the endpoint? Not yet! A function that returns ``I
found a bug'' unconditionally would also report a finding, and it would be
useless. So every acceptance test runs the check twice: once against the
hostile configuration, where it must find the bug, and once against the
boundable one, where it must find nothing. A finding in the first and silence
in the second is evidence about the endpoint. A finding in both is evidence
about the checker, and it means you should go fix the checker before you tell
anybody anything.

\paragraph{Something you can actually deploy, and its ceiling.}
The third part is a proxy you put in front of a real API. It reads every
request strictly, and---this is the important bit---forwards the
\emph{canonical} form rather than the bytes that arrived, so the endpoint
behind it and whatever signed the request cannot be looking at two different
spellings of one document.

Its limits are more interesting than its function, so here is the whole
scoreboard against the three attacks we reproduce:

\begin{center}
\begin{tabular}{lll}
\toprule
Attack & Condition & Through the proxy \\
\midrule
Duplicate-key split & C2 & Rejected, unconditionally \\
Parameter-defaulting bypass & C1 & Rejected \emph{only if} the route names the parameters \\
URI alias collapse & C2 & \textbf{Passes} \\
\bottomrule
\end{tabular}
\end{center}

Look at the third row. That is not a bug we haven't got to yet. Think about
what the proxy actually sees when \texttt{GET /v1/users/me} goes past: a
method, a path, and no body at all. There is nothing in there to canonicalize
and nothing to compare. For the proxy to know that \texttt{me} and
\texttt{usr\_99812} are the same person, it would have to understand the
application's identity model, at which point it is not a proxy any more, it is
a second copy of the application, and now you have two of them to keep in
agreement. The honest statement is that a proxy cannot close this, and no
amount of engineering effort changes that.

So we wrote a test that asserts the attack \emph{succeeds}. It is in the suite,
it is green, and if somebody later makes the README claim that the proxy stops
all three attacks, that test goes red. I recommend this trick generally. The
claims that quietly grow in the retelling are exactly the ones nobody thought
to write a test for, because who writes a test that something doesn't work?

And the second row deserves a moment too. C1 is closed only if you list the
parameters for that route in the configuration. If you don't, the defaulting
bypass sails straight through. That means the proxy's usefulness depends on
somebody doing homework, and the natural way to do that homework is to point
the classifier at the API's specification and read off which parameters it
says are optional. The tool that finds the problem and the tool that fixes it
are the same tool, used twice.

\section{The Survey We Inherited}
\label{sec:survey}

This work inherited a survey. Somebody had gone through 24 production APIs,
classified each against C1--C4 by reading the documentation, and reported that
58.3\% (14 of 24) fail outright, 25.0\% (6 of 24) qualify conditionally, and
16.7\% (4 of 24) qualify. Those percentages are quotable, and they had already
been quoted.

We put the table into a data file so the tool could count it, because we had
a rule that no number gets typed into this paper by hand. The tool counted 4,
7, and 13.

That is not what the prose said. So one of them is wrong, and the useful
question is which.

It's the prose. Square's \texttt{POST /v2/payments} is sitting right there in
the table as a conditional row, and the 25\% figure doesn't include it, and the
58\% figure has swallowed it. The rows were fine all along. Somebody counted
them by hand once, and the count went into the text, and after that everyone
including us quoted the text instead of recounting the rows. An earlier draft
of this very paper repeated the published figures.

I want to dwell on this for a moment, because the arithmetic is trivial and
the lesson isn't. We did not find this by being clever. We found it because we
had made it impossible to write the number down without computing it, and the
computation disagreed. The rule cost almost nothing to adopt and it caught
something on its first application. That is a better argument for the rule
than any amount of methodology talk.

And now the part that the corrected number does not entitle anyone to skip.
Twenty-four endpoints chosen to cover a spread of industries is not a random
sample of anything. It was assembled for coverage, which is the right way to
assemble a set of illustrations and the wrong way to assemble a basis for
percentages. Worse, every one of these classifications came from reading
documentation, and two of our four conditions are not visible in documentation
even in principle. No specification will tell you what a server quietly fills
in from your account settings. No specification will tell you whether an
effect depends on state that moved while you weren't looking. Those are facts
about running code.

So 54\% is not a fact about the world's APIs. Neither was 58\%. What we have
is a taxonomy with two dozen worked examples, which is genuinely useful and is
not a measurement, and the difference between those two things is most of what
this section is about.

Because that distinction is easy to lose in transit, every rendering of this
table carries a banner saying it is a hypothesis---including the generated
\LaTeX{}, where the banner is the caption. You cannot \texttt{\textbackslash
input} the table into a document without the caveat coming along with it. We
did that deliberately, after noticing how easily we had ourselves absorbed the
original percentages.

% ===== PLACEHOLDER: results not yet written up =====
\section{Does It Actually Work?}
\label{sec:results}
% RESULTS PLACEHOLDER — every number below arrives from ../results/, written
% by scripts/regen-results.sh and regenerated in CI. Do not type numbers here.

Every number in this section is generated. None of it is typed. The tables
below are written by \texttt{scripts/regen-results.sh}, the build regenerates
them and fails if they have drifted from the code, and the reason for all that
machinery is in Section~\ref{sec:survey}: we tried the other way and it put a
wrong percentage into circulation.

\subsection{Where decoders disagree}

\begin{center}
% Generated by scripts/regen-results.sh. Do not edit.
\begin{tabular}{llr}
\toprule
Specification & Verdict & Failing conditions \\
\midrule
\texttt{iso20022-pain001} & Qualified & none \\
\texttt{s3-put-if-match} & Qualified & none \\
\texttt{slack-chat-postmessage} & NonQualifying & C1, C2, C3 \\
\texttt{stripe-charges} & Conditional & none \\
\bottomrule
\end{tabular}

\end{center}

\subsection{The survey is a hypothesis}
% The banner is part of the generated table's caption on purpose: this table
% must not be readable without the caveat that produced it.
% Generated by `transit survey --format tex`. Do not edit.
\begin{table}[t]
\caption{HYPOTHESIS, NOT DATA. These 24 classifications are derived from published documentation and have not been measured against a running endpoint. C1 and C4 are properties of an implementation, which no document reveals. The sample was chosen for sector coverage, so it is a convenience sample and the shares below are not a population estimate. Verifying these rows by live probing is work this repository has not done.}
\begin{tabular}{lllcccc l}
\toprule
Sector & Provider & Endpoint & C1 & C2 & C3 & C4 & Hypothesis \\
\midrule
Financial / Payments & ISO 20022 & \texttt{pain.001 credit transfer} & Satisfied & Satisfied & Satisfied & Satisfied & Qualified \\
Financial / Payments & Stripe & \texttt{POST /v1/charges} & ClientDependent & Satisfied & ClientDependent & Fails & Conditional \\
Financial / Payments & Square & \texttt{POST /v2/payments} & ClientDependent & Satisfied & ClientDependent & Fails & Conditional \\
Financial / Payments & Plaid & \texttt{POST /link/token/create} & Fails & Satisfied & Fails & Fails & NonQualifying \\
Cloud Infrastructure & AWS EC2 & \texttt{RunInstances} & Fails & Fails & ClientDependent & Fails & Conditional \\
Cloud Infrastructure & AWS S3 & \texttt{PUT /object} & Fails & Fails & Satisfied & Satisfied & Qualified \\
Cloud Infrastructure & GCP Compute & \texttt{instances.insert} & Fails & Fails & ClientDependent & Fails & Conditional \\
Cloud Infrastructure & Azure ARM & \texttt{PUT /subscriptions/...} & Fails & Fails & Satisfied & Satisfied & Conditional \\
Cloud Infrastructure & Cloudflare & \texttt{POST /dns\_records} & Fails & Fails & Fails & Fails & NonQualifying \\
Developer Tools & GitHub & \texttt{POST /user/repos} & Fails & Fails & Fails & Fails & NonQualifying \\
Developer Tools & GitHub & \texttt{PUT /repos/\{o\}/\{r\}/contents} & Satisfied & Satisfied & Satisfied & Satisfied & Qualified \\
Developer Tools & GitLab & \texttt{POST /api/v4/projects} & Fails & Fails & Fails & Fails & NonQualifying \\
Developer Tools & Docker Engine & \texttt{POST /containers/create} & Fails & Fails & Fails & Satisfied & Conditional \\
SaaS / Productivity & Google Calendar & \texttt{POST /events} & Fails & Fails & Fails & Fails & NonQualifying \\
SaaS / Productivity & Google Workspace & \texttt{POST /gmail/v1/users/me/messages/send} & Fails & Fails & Fails & Fails & NonQualifying \\
SaaS / Productivity & Microsoft Graph & \texttt{POST /v1.0/me/sendMail} & Fails & Fails & Fails & Fails & NonQualifying \\
SaaS / Productivity & Slack & \texttt{POST /chat.postMessage} & Fails & Fails & Fails & Fails & NonQualifying \\
SaaS / Productivity & Jira & \texttt{POST /rest/api/3/issue} & Fails & Fails & Fails & Fails & NonQualifying \\
SaaS / Productivity & Salesforce & \texttt{POST /sobjects/Account} & Fails & Satisfied & Fails & Fails & NonQualifying \\
SaaS / Productivity & HubSpot & \texttt{POST /crm/v3/objects/contacts} & Fails & Fails & Fails & Fails & NonQualifying \\
Identity / Security & Okta & \texttt{POST /api/v1/users} & Fails & Fails & Fails & Fails & NonQualifying \\
Identity / Security & Auth0 & \texttt{POST /api/v2/users} & Fails & Fails & Fails & Fails & NonQualifying \\
Database / Storage & HashiCorp Vault & \texttt{POST /v1/secret/data} & Fails & Fails & Satisfied & Satisfied & Qualified \\
Database / Storage & Pinecone & \texttt{POST /vectors/upsert} & Fails & Fails & Satisfied & Fails & Conditional \\
\bottomrule
\end{tabular}
% Qualified: 4 (16.7%) of a convenience sample of 24, not a population estimate
% Conditional: 7 (29.2%) of a convenience sample of 24, not a population estimate
% NonQualifying: 13 (54.2%) of a convenience sample of 24, not a population estimate
\end{table}


\section{What We Still Don't Know}
\label{sec:openwork}

This is the section I would read first if somebody handed me this paper, so I
have tried to make it the honest one rather than the modest-sounding one.

\paragraph{We mostly cannot check the thing we are asking you to check.}
Start with the worst of it. Two of our four conditions---C1 and C4---are
properties of running code, not of published interfaces. If you do not have
the source, classifying an endpoint means inferring what the implementation
does from what its documentation says and from poking at it, and the entire
credibility of any survey rests on how that inference gets made.

Our classifier answers \texttt{Unknown} whenever a document cannot settle the
question, and refuses to guess. That is honest. It is also not a measurement,
and I do not want the honesty to be mistaken for rigour. The way to turn these
rows into evidence is to probe endpoints that are actually running. We built
the gate that would let somebody do that safely---it demands a written
declaration of authorization, a contact address, and a command-line flag, and
refuses otherwise---and then we did not run it. So the empirical half of this
paper is a plan, not a result, and Section~\ref{sec:survey} is a hypothesis
wearing a table.

\paragraph{Checking somebody who isn't helping you.}
Everything above assumes you may ask the endpoint questions. Suppose you may
not: you are an auditor, or a customer doing diligence, and you have no
account and no permission. Is there any way to establish whether a third-party
API is effect-binding? We do not know of one that is both sound and lawful,
and I suspect this is hard rather than merely unexplored.

\paragraph{What do you do about the ones that can't be fixed?}
Thirteen of the twenty-four endpoints in that survey cannot satisfy these
conditions without changing their public interfaces, and most of them are not
going to. Telling those operators ``adopt C1--C4'' is not advice, it is a
complaint. The real answers are presumably about detecting the damage
afterwards, keeping the blast radius small, and compensating controls---and we
have developed none of that here.

\paragraph{Two calls in a row.}
Now for the one I think is actually the important open problem. This entire
paper is about one authorized message producing one effect. That is not what
agents do. An agent runs a workflow: call something, look at the answer,
decide what to call next, feed the result into the one after that. Each step's
input depends on the previous step's output, so the thing a user approved at
the start was not a message at all---it was closer to a goal.

Does effect binding compose across such a sequence? We don't know. What would
it even mean to authorize a \emph{plan} rather than a message, when the plan's
later steps are not determined until the earlier ones return? We don't know
that either, and I am not sure the vocabulary exists yet.

We have said nothing about it in this paper. It is very likely the question
the whole agent business actually poses, and the four conditions here are, at
best, what you would need to have settled before you could begin asking it.

\bibliographystyle{plain}
\bibliography{references}

\end{document}
